\section{Outro}\label{sec:outro}

This whitepaper has outlined the ``\textbf{Onchain Randomness Problem}'' and why it now affects every AO process: HyperBEAM's runtime no longer offers \texttt{math.random}, because node-local randomness forks process state and can be steered by whoever controls the seed.

RandAO brings it back safely, as the first tool in a suite of cryptographic tools for the new age of the internet. One library gives every process deterministic randomness for free and upgrades the same call to verifiable randomness when value is at stake. Its flagship tier keeps the any-honest time-lock commit-reveal RandAO was founded on: as long as a single validator is honest, nobody can predict or bias a draw, and a validator who withholds can only delay it. Checking each reveal with the puzzle's one-time trapdoor inside the node, through the \texttt{timelock@1.0} device, makes that guarantee fast and cheap. Time-lock encryption follows on the same puzzles.

The \$RNG token pays for randomness: one token to each validator that served a request and one to the protocol. RandAO runs on testnet today, with launch validators operated by RandAO Labs; opening the validator set comes next.
